% !TeX spellcheck = de_DE

Wie oben bereits beschrieben, ist das Ziel dieser Arbeit die Untersuchung, ob die Qualität automatisch generierter Texte in Leichter Sprache Plus (LS+) durch einen metrisch gesteuerten Feedbackloop verbessert werden kann.
Der entwickelte Proof of Concept erweitert die bestehende LS+-Pipeline um eine zusätzliche Qualitätssicherung.
Die generierten Texte werden automatisiert bewertet und bei Bedarf über eine iterative Überarbeitung  strukturiert verbessert.

Die Grundannahme des Ansatzes besteht darin, dass sich Qualitätsmängel in LS+-Texten zumindest teilweise über Metriken erkennen lassen.
Werden Abweichungen von den definierten Qualitätsanforderungen festgestellt, können diese Informationen genutzt werden, um einem Large Language Model gezieltes Feedback für eine weitere Textgenerierung bereitzustellen.
Durch die wiederholte Bewertung und Überarbeitung entsteht ein geschlossener Feedbackkreislau, dessen Ziel die iterative Verbesserung der Textqualität ist.